\documentclass[11pt,a4paper]{article}
\usepackage[utf8]{inputenc}
\usepackage{amsmath,graphicx}
\usepackage[margin=25mm]{geometry}
\usepackage{tikz}
\usetikzlibrary{arrows.meta,positioning}
\newcommand{\dif}{\mathrm{d}}
\newcommand{\Vf}{V_{F}}
\definecolor{myblue}{RGB}{0,80,160}
\title{Il diodo di potenza PIN}
\author{D. Bagnara}
\date{\today}
\begin{document}
\maketitle
\tableofcontents
\begin{abstract}
Nota sul diodo PIN.
\end{abstract}
\section{Struttura}
\label{sec:struct}
La regione $i$ e' descritta in \ref{sec:struct}. Vedi \cite{baliga}.
\begin{equation}\label{eq:vd}
V_{drift}=\frac{W_D^2}{(\mu_n+\mu_p)\tau}
\end{equation}
\section{Reverse recovery}
Come da \eqref{eq:vd} e \cite{lauritzen1991}.
\begin{thebibliography}{9}
\bibitem{baliga} B. J. Baliga, \emph{Fundamentals of Power Semiconductor Devices}, Springer, 2008.
\bibitem{lauritzen1991} P. O. Lauritzen and C. L. Ma, "A simple diode model with reverse recovery," IEEE Trans. Power Electron., 1991.
\end{thebibliography}
\end{document}
